\documentclass[12pt,a4paper]{article}
\usepackage{amsmath,amsfonts,amssymb,mathrsfs,tikz,times,pifont}
\usepackage{enumitem}
\usepackage{multicol}
\newcommand\circitem[1]{%
\tikz[baseline=(char.base)]{
\node[circle,draw=gray, fill=red!55,
minimum size=1.2em,inner sep=0] (char) {#1};}}
\newcommand\boxitem[1]{%
\tikz[baseline=(char.base)]{
\node[fill=cyan,
minimum size=1.2em,inner sep=0] (char) {#1};}}
\setlist[enumerate,1]{label=\protect\circitem{\arabic*}}
\setlist[enumerate,2]{label=\protect\boxitem{\alph*}}
%%%::::::by chnini ameur :::::::%%%
\everymath{\displaystyle}
\usepackage[left=1cm,right=1cm,top=1cm,bottom=1.7cm]{geometry}
\usepackage{array,multirow}
\usepackage[most]{tcolorbox}
\usepackage{varwidth}
\usepackage{mathrsfs} % pour \mathscr
\usepackage{calligra}
\newcommand{\Dcalligra}{\text{\calligra D}}
\tcbuselibrary{skins,hooks}
\usetikzlibrary{patterns}
%%%::::::by chnini ameur :::::::%%%
\newtcolorbox{exa}[2][]{enhanced,breakable,before skip=2mm,after skip=5mm,
colback=yellow!20!white,colframe=black!20!blue,boxrule=0.5mm,
attach boxed title to top left ={xshift=0.6cm,yshift*=1mm-\tcboxedtitleheight},
fonttitle=\bfseries,
title={#2},#1,
% varwidth boxed title*=-3cm,
boxed title style={frame code={
\path[fill=tcbcolback!30!black]
([yshift=-1mm,xshift=-1mm]frame.north west)
arc[start angle=0,end angle=180,radius=1mm]
([yshift=-1mm,xshift=1mm]frame.north east)
arc[start angle=180,end angle=0,radius=1mm];
\path[left color=tcbcolback!60!black,right color = tcbcolback!60!black,
middle color = tcbcolback!80!black]
([xshift=-2mm]frame.north west) -- ([xshift=2mm]frame.north east)
[rounded corners=1mm]-- ([xshift=1mm,yshift=-1mm]frame.north east)
-- (frame.south east) -- (frame.south west)
-- ([xshift=-1mm,yshift=-1mm]frame.north west)
[sharp corners]-- cycle;
},interior engine=empty,
},interior style={top color=yellow!5}}
%%%%%%%%%%%%%%%%%%%%%%%

\usepackage{fancyhdr}
\usepackage{eso-pic}         % Pour ajouter des éléments en arrière-plan
% Commande pour ajouter du texte en arrière-plan
%\AddToShipoutPicture{
%    \AtTextCenter{%
%        \makebox[0pt]{\rotatebox{80}{\textcolor[gray]{0.7}{\fontsize{5cm}{5cm}\selectfont PGB}}}
%    }
%}
\usepackage{lastpage}
\fancyhf{}
\pagestyle{fancy}
\renewcommand{\footrulewidth}{1pt}
\renewcommand{\headrulewidth}{0pt}
\renewcommand{\footruleskip}{10pt}
\fancyfoot[R]{
\color{blue}\ding{45}\ \textbf{2026}
}
\fancyfoot[L]{
\color{blue}\ding{45}\ \textbf{PGB}
}
\cfoot{\bf
\thepage /
\pageref{LastPage}}
\begin{document}
\renewcommand{\arraystretch}{1.5}
\renewcommand{\arrayrulewidth}{1.2pt}
\begin{tikzpicture}[overlay,remember picture]
\node[draw=blue,line width=1.2pt,fill=purple,text=blue,inner sep=3mm,rounded corners,pattern=dots]at ([yshift=-2.5cm]current page.north) {\begingroup\setlength{\fboxsep}{0pt}\colorbox{white}{\begin{tabular}{|*1{>{\centering \arraybackslash}p{0.28\textwidth}} |*2{>{\centering \arraybackslash}p{0.2\textwidth}|} *1{>{\centering \arraybackslash}p{0.19\textwidth}|} }
\hline
\multicolumn{3}{|c|}{$\diamond$$\diamond$$\diamond$\ \textbf{Lycée Dindéfélo}\ $\diamond$$\diamond$$\diamond$ }& \textbf{A.S. : 2025/2026} \\ \hline
\textbf{Matière: Mathématiques}& \textbf{Niveau : 4}\textbf{ème} &\textbf{Date: 18/05/2026} & \textbf{Durée : 2 heures} \\ \hline
\multicolumn{4}{|c|}{\parbox[c]{10cm}{\begin{center}
\textbf{{\Large\sffamily Devoir 1 Du 2$ ^\text{\bf nd} $ Semestre}}
\end{center}}} \\ \hline
\end{tabular}}\endgroup};
\end{tikzpicture}
\vspace{3cm}

\section*{\underline{Exercice 1 :} Droites remarquables du triangle (6 points)}

Compléter les phrases suivantes en utilisant les termes : \textbf{médiatrices}, \textbf{hauteurs}, \textbf{centre de gravité}, \textbf{orthocentre}, \textbf{bissectrices} ou \textbf{médianes}.

\begin{enumerate}
    \item Les trois \dotfill d'un triangle se coupent au centre du cercle inscrit. \hfill \textbf{(1,5 pt)}
    \item Le point de rencontre des trois \dotfill est appelé \textbf{orthocentre}. \hfill \textbf{(1,5 pt)}
    \item Le \dotfill est le point d'intersection des trois médianes. \hfill \textbf{(1,5 pt)}
    \item Les trois \dotfill d'un triangle se coupent au centre du cercle circonscrit. \hfill \textbf{(1,5 pt)}
\end{enumerate}

\section*{\underline{Exercice 2 :} Identités remarquables et Factorisation (7 points)}

\begin{enumerate}
    \item En utilisant les identités remarquables $(a+b)^2$, $(a-b)^2$ ou $(a-b)(a+b)$, développer et réduire les expressions suivantes : \textbf{(3 pts)}

\begin{multicols}{3}
\setlength{\columnseprule}{0.1mm}
\begin{enumerate}
\item[]
\(
\begin{aligned}
A&=(x + 5)^2\\
&=\\
&=\\
\end{aligned}
\)
\item[]
\(
\begin{aligned}
B&=(3x - 2)^2\\
&=\\
&=\\
\end{aligned}
\)
\item[]
\(
\begin{aligned}
C&=(x - 4)(x + 4)\\
&=\\
&=\\
\end{aligned}
\)
\end{enumerate}
\end{multicols} 
    
\item Développer et réduire les expressions suivantes : \textbf{(2 pts)}
 \begin{multicols}{2}
\setlength{\columnseprule}{0.1mm}
\begin{enumerate}
\item[]
\(
\begin{aligned}
D&=3(2x - 5) + 4x(x + 2)\\
&=\\
&=\\
&=\\
\end{aligned}
\)
\item[]
\(
\begin{aligned}
E&=(2x + 3)(x - 4) - (x^2 - 5)\\
&=\\
&=\\
&=\\
\end{aligned}
\)
\end{enumerate}
\end{multicols}

    \item Factoriser les expressions suivantes en recherchant un facteur commun : \textbf{(2 pts)}

\begin{multicols}{2}
\setlength{\columnseprule}{0.1mm}
\begin{enumerate}
\item[]
\(
\begin{aligned}
F&=(x + 3)(2x - 1) + (x + 3)(x + 5)\\
&=\\
&=\\
&=\\
\end{aligned}
\)
\item[]
\(
\begin{aligned}
G&=(x - 2)^2 - 3(x - 2)\\
&=\\
&=\\
&=\\
\end{aligned}
\)
\end{enumerate}
\end{multicols}     
    
\end{enumerate}
\newpage

\section*{\underline{Exercice 3 :} Équations et Inéquations (7 points)}
\begin{enumerate}
    \item Résoudre dans $\mathbb{Q}$ les équations suivantes : \textbf{(4 pts)}

\begin{multicols}{3}
\setlength{\columnseprule}{0.1mm}
\begin{enumerate}
\item[]
\(
\begin{aligned}
5x - 8 &= 2x + 7\\
... &= 7 + 8\\
3x &=... \\
... &=... \\
S=\{ \quad\quad \}
\end{aligned}\\
\)
\item[]
\(
\begin{aligned}
3(x - 2) &= 4x + 5\\
... &=... \\
... &= 5 + 6\\
... &= ...\\ 
x&=... \\
S=\{ \quad\quad \}
\end{aligned}
\)
\item[]
\(
\begin{aligned}
(2x - 4)(3x + 9) &= 0\\
2x - 4 = 0 &\textbf{ ou } .....= 0\\
.....= 4 &\textbf{ ou } 3x =..... \\
.....= ... &\textbf{ ou }..... = .....\\
.....= ... &\textbf{ ou }...... = -3\\
S=\{ \quad;\quad \}
\end{aligned}
\)
\end{enumerate}
\end{multicols}    
    
    \item Résoudre dans $\mathbb{Q}$ les inéquations suivantes et représenter les solutions sur une droite graduée : \textbf{(3 pts)}
\begin{multicols}{2}
\setlength{\columnseprule}{0.1mm}
\begin{enumerate}
\item[]
\(
\begin{aligned}
2x + 5 &< 13\\
.... &< 13 - ...\\
... &< ...\\
x &< ...\\
S=
\end{aligned}
\)
\item[]
\(
\begin{aligned}
-3x + 4 &\geq 10\\
... &\geq 1....\\
... &\geq ...\\
...\quad &...\quad... \\
...\quad &...\quad...\\
S=
\end{aligned}
\)

\end{enumerate}
\end{multicols}  
\end{enumerate}
\end{document}